\documentclass[a4paper,11pt]{article}
\usepackage{exam-style-842-unit}
\examinehead{842 信息化概论 · 第 8 单元练习册}{练习二\quad 移动电商、支付物流与电子政务}
\begin{document}

\unitcover{842 信息化概论 · 单元练习册}{第 8 单元\quad 电子商务与电子政务}{练习二\quad 移动电商、支付物流与电子政务（共 100 分）}{约 70 分钟}{移动电子商务的概念、特点与运营模式；电子支付与常用支付工具；网络银行与互联网金融模式；电子商务物流与供应链管理；电子政务的概念、4 种应用模式、电子政务服务体系与电子公务。}

\begin{knowmap}
\textbf{移动电子商务：}移动终端＋网络＋交易活动＝移动电子商务；主要特点是灵活、简单、方便，可按消费者个性化需求定制\\[2pt]
\textbf{运营模式（按参与角色）：}内容提供商 · 移动门户 · 移动运营商 · WAP 网关提供商 · 移动直接面向客户 · 移动虚拟社区\\[2pt]
\textbf{支付链：}支付发起 $\rightarrow$ 支付指令交换与清算（含轧差）$\rightarrow$ 结算；电子支付类型按发起方式分为网上、电话、移动、POS、ATM 等\\[2pt]
\textbf{支付工具：}电子现金（数字现金）· 银行卡 · 电子支票 · 电子钱包 · 网络银行（“3A 银行”）$\Rightarrow$ 安全靠服务器安全、身份识别（数字证书/CA）、SSL 协议与多种认证介质\\[2pt]
\textbf{物流与供应链：}电子商务物流向信息化、自动化、智能化发展；供应链管理实现信息网状传递、单证标准化、对消费者需求即时反应\\[2pt]
\textbf{电子政务：}政府机关利用网络计算机及通信等信息技术优化管理与服务；4 种模式 G2G、G2C、G2B、G2E；服务体系＝公共服务前台（门户网站＋行政服务中心）＋政府办公后台（OA＋资源共享＋协同办公）；电子公务 eGBCP
\end{knowmap}

\begin{keyterms}
\textbf{移动电子商务}：基于无线网络，运用手机、掌上电脑等移动通信设备与因特网有机结合所进行的电子商务活动。\\
\textbf{WAP}：把 Internet 技术与无线网络技术结合、针对移动通信设备接入互联网而设计的通信协议规范。\\
\textbf{WAP 网关}：对 Web 服务器信息进行转换、使移动终端能够访问 Web 内容的网关设备。\\
\textbf{移动虚拟社区}：为有相同爱好、经历或业务相关的网络用户提供交流与分享场所的虚拟社区。\\
\textbf{电子支付}：单位或个人直接或授权他人通过电子终端发出支付指令，实现货币支付与资金转移的行为。\\
\textbf{清算}：结算之前对支付指令进行发送、对账、确认的处理，还可能包括轧差。\\
\textbf{结算}：双方或多方对支付交易相关债务的清偿。\\
\textbf{电子现金}：又称数字现金，是以数据形式流通、模拟使用现金的加密序列数，具有货币价值、可交换性、可存储性和不可重复性。\\
\textbf{电子支票}：借鉴纸张支票转移支付的优点，利用数字传递将钱款从一个账户转移到另一个账户的电子付款形式。\\
\textbf{电子钱包}：电子商务活动中顾客常用的小额支付工具，用于携带信用卡、借记卡等支付工具。\\
\textbf{网络银行}：又称网上银行、在线银行，银行利用 Internet 向客户提供开户、查询、转账、理财等服务的虚拟银行柜台。\\
\textbf{3A 银行}：网络银行的别称，可在任何时间（Anytime）、任何地点（Anywhere）、以任何方式（Anyway）提供金融服务。\\
\textbf{第三方支付}：具备一定实力和信誉保障的非银行机构，借助信息技术与各大银行签约，在用户与银行支付结算系统间建立连接的电子支付模式。\\
\textbf{电子政务}：政府机关利用网络计算机及通信等信息技术，实现管理和服务无缝集成、业务流程优化，并通过政府网站提供一体化管理和服务。\\
\textbf{电子政务服务体系}：依托国家电子政务建设平台，以公共服务为中心、以满足公众需求为导向，对内实现信息共享与协同办公、对外提供无缝隙“一站式”服务的一体化服务体系。\\
\textbf{电子公务}：利用信息化手段，高效实现政府、企业、公众（社区）在生产公共产品过程中的事务，其模式为 eGBCP。
\end{keyterms}

\begin{contrast}{电子现金 / 电子支票 / 电子钱包（每行依次为：定义·主要特点·典型用途）}
\citem{电子现金}{以数据形式流通、模拟使用现金的加密序列数}{具有货币价值、可交换性、可存储性、不可重复性}{小额网上支付，要求匿名性与离线支付}
\citem{电子支票}{利用数字传递将钱款从一个账户转移到另一个账户的电子付款形式}{与传统支票工作方式相同，加密后比数字现金更易流通，可与 EDI 结合}{企业与企业之间的大额付款}
\citem{电子钱包}{携带信用卡、借记卡等支付工具的软件}{实行密码管理，需申请电子安全证书}{小额购物或购买小商品}
\end{contrast}

\begin{contrast}{G2G / G2C / G2B / G2E（每行依次为：主体关系·主要内容·典型应用）}
\citem{G2G}{政府与政府、政府不同机构和部门之间}{电子公文处理、电子政策法规系统、电子财政管理、政府网络业绩管理}{跨部门公文流转、财政预算与执行情况报告}
\citem{G2C}{政府与公众之间}{教育培训、公众福利、电子就业、电子税务、心理咨询、电子医疗、电子证件服务}{网上办证、网上报税、网上就业服务}
\citem{G2B}{政府与企业之间}{信息咨询、电子采购与招标、电子税务、电子证件办理、中小企业服务}{政府采购招标、企业网上报税}
\citem{G2E}{政府与政府雇员之间}{电子公文处理、电子政策法规、电子培训、电子财务管理、公务员业绩评价}{内网办公自动化、公务员在线培训与考核}
\end{contrast}

\begin{contrast}{电子政务 / 电子公务（每行依次为：强调角色·参与主体·范畴侧重）}
\citem{电子政务}{强调政府的管理角色}{以政府为中心}{政府与非营利组织、营利组织、公民、政府之间的行政事务}
\citem{电子公务}{强调政府的服务角色}{政府、企业、公众（社区）共同参与}{生产公共产品过程中的公共事务，含非营利组织、企业、市民参与的活动}
\end{contrast}

\newpage
\begin{contrast}{公共服务前台 / 政府办公后台（每行依次为：组成·服务对象·主要职能）}
\citem{公共服务前台}{政府门户网站＋行政服务中心}{社会公众}{信息发布、政务公开、公众参与、在线办公与“一站式”行政审批服务}
\citem{政府办公后台}{办公自动化系统＋政务信息资源共享平台＋协同办公集成平台}{政府内部部门与公务员}{办公自动化、打破“信息孤岛”实现信息共享、跨部门协同办公}
\end{contrast}

\qtypename{一、单项选择题}{10 题，每题 3 分，共 30 分}
\q 移动电子商务可以概括为（\quad）
\choices{A. 移动终端＋网络＋交易活动}{B. PC＋宽带＋交易活动}
{C. 电话＋传真＋交易活动}{D. 物流＋资金流＋信息流}

\q 移动电子商务的主要特点是（\quad）
\choicesline{A. 灵活、简单、方便}{B. 固定、复杂、昂贵}{C. 只能在室内使用}{D. 必须使用现金支付}

\q 电子支付按支付指令发起方式分类，不包括（\quad）
\choicesline{A. 网上支付}{B. 电话支付}{C. 移动支付}{D. 现金支付}

\q 电子现金必须防止复制和重复使用，这体现其（\quad）属性。
\choicesline{A. 货币价值}{B. 可交换性}{C. 可存储性}{D. 不可重复性}

\q 不受时间、空间限制，能在任何时间、任何地点、以任何方式提供金融服务的网络银行又被称为（\quad）
\choicesline{A. 3A 银行}{B. 3C 银行}{C. 5C 银行}{D. 3P 银行}

\q 目前用信用卡进行网上支付的国际标准是（\quad）
\choicesline{A. SET 协议}{B. FTP 协议}{C. SMTP 协议}{D. TCP 协议}

\q 第三方支付狭义上是指（\quad）
\choices{A. 商业银行提供的网上银行服务}{B. 非银行机构在用户与银行支付结算系统间建立连接的电子支付模式}
{C. 消费者之间的直接转账}{D. 政府财政支付}

\q 电子政务应用模式中，政府与政府之间、政府不同机构和部门之间的是（\quad）
\choicesline{A. G2G}{B. G2C}{C. G2B}{D. G2E}

\q 电子政务服务体系的灵魂和价值在于（\quad）
\choicesline{A. 服务}{B. 管理}{C. 监管}{D. 审批}

\q 电子公务的 eGBCP 模式中，P 代表（\quad）
\choicesline{A. 公众}{B. 企业}{C. 公共设施与公共环境}{D. 政府}

\qtypename{二、多项选择题}{3 题，每题 4 分，共 12 分；每题有两个或两个以上正确项}
\q 移动电子商务的运营模式（按参与角色划分）包括（\quad）
\choicesline{A. 内容提供商}{B. 移动门户}{C. 移动运营商}{D. WAP 网关提供商}

\q 网络银行安全体系通常包括（\quad）
\choices{A. 服务器安全（防火墙、安全服务器、24 小时监控）}{B. 身份识别（数字证书、数字签名等）}
{C. 安全协议（SSL 等）}{D. 认证介质（动态口令卡、U 盾等）}

\q 电子政务的应用模式包括（\quad）
\choicesline{A. G2G}{B. G2C}{C. G2B}{D. G2E}

\qtypename{三、判断题}{6 题，每题 2 分，共 12 分；正确的画“√”，错误的画“×”并写出改正后的表述}
\q 移动电子商务是电子商务的一个子集和延伸。\quad（\quad）\quad 改正：\fixblank

\q 电子现金可以复制和重复使用，不会带来安全问题。\quad（\quad）\quad 改正：\fixblank

\q SSL 协议由 Netscape 研制，用于在两个通信方之间提供秘密而可靠的连接。\quad（\quad）\quad 改正：\fixblank

\q 电子政务就是电子公务，二者完全等同。\quad（\quad）\quad 改正：\fixblank

\q 第三方支付机构本身就是银行，可以吸收公众存款。\quad（\quad）\quad 改正：\fixblank

\q 我国电子政务服务体系由公共服务前台和政府办公后台两大平台构成。\quad（\quad）\quad 改正：\fixblank

\qtypename{四、填空题}{6 空，每空 1 分，共 6 分}
\q 移动电子商务＝移动终端＋\fillblanks[1.6cm]＋交易活动。

\q 网络银行又称“3A 银行”，3A 指 Anytime、Anywhere 和\fillblanks[2.0cm]。

\q 国际通行的网上支付工具主要有电子现金、\fillblanks[2.4cm]、电子支票、电子钱包和网络银行等。

\q 电子政务根据政府、企业和社会公众 3 个主体之间的关系分为 G2G、G2C、\fillblanks[1.8cm]和 G2E 4 种模式。

\q 我国电子政务服务体系由\fillblanks[2.8cm]（政府门户网站和行政服务中心）与政府办公后台两大平台构成。

\q eGBCP 模式中，G 代表政府、B 代表企业、C 代表公众和社区、P 代表\fillblanks[2.8cm]。

\qtypename{五、简答题}{4 题，每题 6 分，共 24 分}
\q 什么是移动电子商务？简述其主要特点。\anslines{2.4cm}

\q 简述移动电子商务的运营模式（按参与角色划分）。\anslines{2.4cm}

\q 简述电子支付的含义、主要类型和常用支付工具。\anslines{2.4cm}

\q 简述电子政务的概念与 4 种应用模式。\anslines{2.4cm}

\newpage
\qtypename{六、综合题}{1 题，共 16 分}
\q 某市推进“一网通办”：政府门户网站和行政服务中心统一受理事项，后台由办公自动化系统、政务信息资源共享平台和协同办公集成平台支撑；市民可用手机 App 办理社保、公积金等业务，并通过移动支付缴费；同时该市把本地农产品放到移动电商平台销售，依托电商物流配送。请回答：
（1）结合材料说明电子政务服务体系的两大平台及其作用；
（2）区分电子政务与电子公务，并指出材料中哪些属于电子公务的范畴；
（3）说明移动支付在“一网通办”和农产品移动电商中分别起什么作用；
（4）说明电子商务物流与供应链管理对农产品上行的意义。\anslines{5cm}

\end{document}
